\subsection{A frontier open-weight model, served through a relay, reproduces the reversal and collapses on the cross-hospital layer}
\label{sec:results-frontier}

The local matrix establishes the four reversals for open-weight checkpoints. A
natural objection is that they are a property of models small enough to fit on one
GPU. We therefore ran the same measurement stack against a frontier-class
\emph{open-weight} model (DeepSeek-V4.1-Flash~\citep{deepseek2026v41flash}), served through an
OpenAI-compatible relay, without changing any part of the protocol that defines
the object being measured. The relay is a transport, not a different model class:
the model is released as open weights, but the relay is not transparent---we
observe the model id and the served outputs, not the serving stack---so this arm
characterizes the endpoint rather than a local build of the checkpoint.

{\sloppy
\paragraph{Protocol.} 9{,}400 calls. Images are produced by the same operators
(\texttt{shifts.apply\_shift}, \texttt{make\_blank}), re-encoded at a 1{,}024-pixel
long edge (JPEG $q{=}95$) and recorded with their transmitted geometry; prompts,
label sets and the JSON parser are byte-identical to the local arm. The CXR panel
is NIH ($400$ studies $\times$ \{\texttt{clean}, contrast-severe, noise-severe,
resolution-severe\} in v0, plus $400\times$ two blank contracts), PadChest
($265$ disease-positive $+235$ negative studies, native and blank) and Open-i
($300$ studies, report and blank). The cross-modal panel is 900 images sampled at
patient or lesion level: 300 each from brain MRI, retinal OCT and dermoscopy, six
cells per image.\par}

\paragraph{Two measurement facts about the relay.} First, a 1{,}024-token budget
returns \texttt{choices: []} with \texttt{completion\_tokens}${=}1{,}024$: the
model spends the whole budget on a hidden reasoning trace. The arm uses 8{,}192
and records \texttt{finish\_reason} per cell; five CXR cells still exhausted that
budget and are counted as contract failures rather than dropped. Second, the relay
does not return \texttt{reasoning\_content}, so the output-mode comparison of
Section~\ref{sec:results-schema} cannot be repeated here; we report the contract
axis only.

\begin{table}[t]
\centering
\small
\caption{Frontier model on the CXR and cross-modal panels. ``Distinct'' counts
different status vectors over the dataset's clean/native cells; ``modal'' is the
share of the most common vector. $\Delta_{\text{null}}$ is
$D(\text{real},\text{blank})-D(\text{real},\text{real}')$ with a 95\% patient- or
lesion-clustered bootstrap interval. Dashes: condition not run for that modality.
Open-i is a free-text report layer without structured status vectors or shift
cells, so only its $\Delta_{\text{null}}$ is defined.}
\label{tab:api-null}
\begin{tabular}{lrrrrrr}
\toprule
Dataset & $n$ & Distinct & Modal & P1 contrast & P1 noise & $\Delta_{\text{null}}$ [95\% CI] \\
\midrule
NIH        & 400 & 92 & 0.370 & 0.590 & 0.733 & 0.254 [0.212, 0.297] \\
PadChest   & 500 & \textbf{3} & \textbf{0.996} & --- & --- & \textbf{0.014} [0.004, 0.026] \\
Dermoscopy & 300 & 94 & 0.117 & 0.417 & 0.540 & 0.155 [0.093, 0.215] \\
Brain MRI  & 300 & 10 & 0.207 & 0.170 & 0.327 & 0.203 [0.122, 0.284] \\
Retina OCT & 300 & 16 & 0.307 & 0.187 & 0.207 & 0.390 [0.320, 0.460] \\
Open-i     & 300 & --- & --- & --- & --- & 0.283 [0.267, 0.298] \\
\bottomrule
\end{tabular}
\end{table}

Three results follow.

\emph{The responsiveness collapse is not a capacity artefact.} For 500 PadChest
studies the frontier model emits \textbf{three} distinct status vectors, with one
vector covering 99.6\% of them. On 400 NIH studies it emits 92, and on 300
dermoscopy images 94. A model with orders of magnitude more capacity than the
local checkpoints behaves exactly as they do on the cross-hospital layer, and
behaves differently on the in-domain layer. Responsiveness is a property of the
dataset--model pair, not of parameter count.

\emph{A positive null separation is necessary but not sufficient.} Every panel
has $\Delta_{\text{null}}>0$ with an interval excluding zero, which is the
condition P2 asks for. Yet PadChest's separation is 0.014, against 0.254 on NIH
and 0.390 on OCT. On PadChest the model's real answers are all but identical to
each other, so the blank differs from them by only a hair. Reporting
$\Delta_{\text{null}}$ without the responsiveness check would therefore certify a
collapsed model as ``image-dependent''. This is the same ordering the local
matrix produces, now at frontier scale.

\emph{The operating-point reversal travels to a relay-served open-weight model.} On NIH, contrast
degradation raises recall from 0.170 to 0.208 while also raising the positive-call
rate on negative studies from 0.116 to 0.163; resolution raises that rate to
0.204 while recall falls. A report that showed only recall would describe the
contrast cell as an improvement. The direction and the mechanism match
Reversal~3 exactly.

\begin{table}[t]
\centering
\small
\caption{Operating point and contract on the frontier model. ``Positive rate''
is the share of label-negative studies called positive; ``recall'' is the share of
label-positive studies called positive. PadChest's zero row is a collapsed output,
not a conservative one: its three distinct vectors are almost all negative.}
\label{tab:api-op}
\begin{tabular}{llrrrrr}
\toprule
Dataset & Condition & $n_{\text{pos}}$ & $n_{\text{neg}}$ & Positive rate & Recall & Parse \\
\midrule
NIH        & clean             & 106 & 294 & 0.116 & 0.170 & 0.998 \\
NIH        & contrast-severe   & 106 & 294 & 0.163 & 0.208 & 1.000 \\
NIH        & noise-severe      & 106 & 294 & 0.133 & 0.104 & 1.000 \\
NIH        & resolution-severe & 106 & 294 & 0.204 & 0.123 & 1.000 \\
PadChest   & native            & 117 & 383 & 0.000 & 0.000 & 1.000 \\
Dermoscopy & clean             & 128 & 172 & 0.488 & 0.836 & 1.000 \\
Dermoscopy & noise-severe      & 128 & 172 & 0.314 & 0.586 & 1.000 \\
Retina OCT & clean             & 230 &  70 & 0.000 & 0.748 & 1.000 \\
\bottomrule
\end{tabular}
\end{table}

\paragraph{What the cross-modal panel adds.} Responsiveness separates by modality
rather than by difficulty: P1 contrast/noise is 0.59--0.73 on CXR and 0.42--0.54
on dermoscopy, but 0.17--0.33 on brain MRI and 0.19--0.21 on retinal OCT. Dermoscopy is the
non-radiology, surface-reflectance modality and it is \emph{not} the least
responsive; the two grayscale tomographic modalities are. That ordering is the opposite of a ``hard modality'' explanation
and is consistent with the corpus-prior account: the model is least responsive
where its training distribution is thinnest, not where the image is hardest.
The operating point reads the same way from the other side: retinal OCT is the
only panel with a zero positive-call rate on its 70 normal studies while
recalling 0.75 of 230 pathological ones, whereas a single contrast step on NIH
raises the negative-call rate by half again (0.116 to 0.163). A model can be
conservative on one modality and unstable on another; neither number travels
without the other.

\paragraph{Limits of this arm.} One model, accessed through a relay, not a family; a probe grid of
four to six cells per image rather than the 22-cell grid; $2{,}085$ cells had to be
re-run after the provider's quota reset, and the re-run is recorded in the
provenance; five CXR cells exhausted an 8{,}192-token budget. None of these affects
the direction of the three results above, but they bound how far the magnitudes
generalise.
